
\subsubsection{3.1 Experimental Design}

The methodology follows fixed-token experimental design: for each curation parameter, models are trained on identical token budgets, architectures, and hyperparameters, varying only the parameter under study.

\textbf{Independent Variables:}

\begin{itemize}
\item Perplexity Threshold: KenLM 5-gram perplexity percentile cutoff. Levels: p0 (none), p10, p20, p30, p40, p50, p60, p70, p80, p90.
\item Deduplication Stringency: MinHash/exact deduplication configuration. Levels: none, fuzzy\_0.7, fuzzy\_0.85, exact.
\end{itemize}

\textbf{Dependent Variables:}

\begin{itemize}
\item Primary: Benchmark ensemble score computed as mean accuracy across HellaSwag, ARC-Easy, PIQA, and WinoGrande.
\item Secondary: Training loss curves, convergence metrics.
\end{itemize}

\textbf{Controlled Variables:}

\begin{itemize}
\item Model architecture: GPT-2 (125M, 350M, 1B variants)
\item Training hyperparameters: learning rate 6e-4, batch size 512, warmup 2000 steps
\item Random seeds: 3 seeds per configuration
\end{itemize}

\subsubsection{3.2 Dose-Response Analysis}

Polynomial models of increasing order are fit and selected using AIC/BIC:

y = $\beta _0 + \beta _1x + \beta _2x^2$ + ... + $\beta$ₖxᵏ + $\epsilon$

Model selection criteria:
\begin{itemize}
\item If quadratic/cubic significantly outperforms linear: evidence for non-monotonic relationship
\item If higher-order model shows interior peak: optimal threshold exists within parameter range
\end{itemize}

For quadratic models, the optimal threshold is x* = -$\beta _1$ / ($2\beta _{2})$, with 95\% confidence intervals computed via bootstrap resampling (1000 iterations).

\subsubsection{3.3 Mechanism Verification}

Convergence dynamics are tracked via:

\begin{enumerate}
\item \textbf{Convergence AUC:} Area under the loss curve, capturing total training cost.
\item \textbf{Steps to Threshold:} Number of steps to reach target loss value.
\end{enumerate}

Prediction: If noise dilutes learning signal, unfiltered training (p0) should show higher AUC than moderate filtering.

\subsubsection{3.4 Scale Transfer Analysis}

Scale transfer is tested by:
\begin{enumerate}
\item Identifying optimal threshold at 125M scale
\item Validating at 1B scale with optimal threshold
\item Computing transfer ratio: (improvement at 1B) / (improvement at 125M)
\end{enumerate}

Success criterion: Transfer ratio > 0.80 indicates optima are scale-stable within 20\%.
